\chapter{Confusion}

\section*{Definition}
Altered mental state characterized by reduced awareness of surroundings, disorganized thinking, and impaired attention.

\section*{When to See a Doctor}
See your doctor if:
\begin{itemize}
  \item Sudden, worsening, or unexplained confusion requires prompt medical assessment; call emergency services for severe drowsiness, weakness, speech trouble, seizure, or inability to stay awake.
\end{itemize}

\section*{How It Develops}
Acute confusion may reflect delirium from infection, medicines, metabolic disturbance, intoxication, stroke, seizure, or another urgent problem. Slowly progressive changes may reflect dementia or other chronic conditions, but the pattern must be assessed clinically.

\section*{Causes}
\begin{itemize}
  \item Infection/sepsis
  \item Metabolic disorders (electrolyte imbalances)
  \item Medication effects/toxicity
  \item Stroke/TIA
  \item Head injury
\end{itemize}

\section*{Related Symptoms}
\begin{itemize}
  \item Fever
  \item Headache
  \item Weakness
  \item Seizures
  \item Memory loss
\end{itemize}


\section*{Medical Evaluation}
\begin{itemize}
  \item Detailed neurological history includes onset pattern, progression, triggers, associated symptoms, and functional impact.
  \item Comprehensive neurological examination assesses mental status, cranial nerves, motor/sensory function, reflexes, coordination, and gait.
  \item Neuroimaging (CT or MRI) is selected according to the onset, examination, trauma history, and concern for stroke, bleeding, or another structural cause.
  \item Specialized testing such as EEG or lumbar puncture is guided by clinical suspicion; nerve-conduction studies do not usually evaluate acute confusion.
\end{itemize}

\section*{Treatment Options}
\begin{itemize}
  \item Treat the underlying cause, such as hypoglycaemia, infection, medication toxicity, hypoxia, or stroke, and correct urgent metabolic abnormalities.
  \item Provide reassurance, orientation cues, sensory aids, hydration when appropriate, and a calm environment while investigating the cause.
  \item Review medications and substances; involve neurology, geriatrics, psychiatry, or other specialists according to the findings.
\end{itemize}

\section*{Self-Care and Home Management}
\begin{itemize}
  \item Keep a calm, familiar environment with minimal distractions during episodes of confusion.
  \item Family members should gently reorient you to person, place, and time if you become confused.
  \item Ensure someone stays with you during episodes if you are at risk of wandering or falls.
  \item Write down important information (phone numbers, medications, allergies) and keep it with you.
  \item Maintain a regular routine for meals, sleep, and activities to reduce disorientation.
  \item Seek medical evaluation for sudden confusion, which can signal a serious medical problem.
\end{itemize}

\section*{References}
\begin{thebibliography}{99}
\bibitem{confusion:ref1} Marcantonio ER. ``Delirium in hospitalized older adults.'' N Engl J Med. 2017;377(15):1456-1466.
\bibitem{confusion:ref2} Inouye SK, Westendorp RGJ, Saczynski JS. ``Delirium in elderly people.'' Lancet. 2014;383(9920):911-922.
\bibitem{confusion:ref3} Fong TG, Tulebaev SR, Inouye SK. ``Delirium in elderly adults: diagnosis, prevention and treatment.'' Nat Rev Neurol. 2009;5(4):210-220.
\bibitem{confusion:ref4} American Psychiatric Association. ``Diagnostic and Statistical Manual of Mental Disorders.'' 5th ed, Text Revision. APA; 2022.
\bibitem{confusion:ref5} Francis J, Young GB. ``Diagnosis of delirium and confusional states.'' In: UpToDate. Wolters Kluwer; 2025.
\bibitem{confusion:ref6} Cummings JL. ``Alzheimer's disease and other dementias.'' In: Harrison's Principles of Internal Medicine. 21st ed. McGraw-Hill; 2022.
\end{thebibliography}
